\chapter{Connect your accounts}\label{ch:accounts}

\emph{You want to see your organization's projects and start agents that work for you.} For that,
Agern needs two things from you: a Claude credential that pays for your agents, and a GitHub token
that decides which repositories you see and that your agents push with. Both are yours alone:
Agern encrypts them, never shows them again (not to you, not to admins), and uses them only for
agents you start. Every person who uses Agern does this once.

Open \ui{Settings} in the top bar. As a new user, every card says what is missing
(figure~\ref{fig:settings-fresh}).

\screenshot{24-settings-fresh-desktop}{Settings of a new user: no Claude credential and no GitHub token yet.}{settings-fresh}

\section{Your Claude credential}\label{sec:claude-credential}

Agents are the \texttt{claude} command-line tool running without a terminal, and they need a
credential that bills to you. There are two kinds:

\begin{itemize}
  \item an \textbf{OAuth token} from \cmd{claude setup-token}, for a Claude subscription (it starts
    with \texttt{sk-ant-oat01-}), or
  \item an \textbf{Anthropic API key} for pay-as-you-go (it starts with \texttt{sk-ant-api03-}).
\end{itemize}

\begin{enumerate}
  \item On \emph{your own} computer, with Claude Code installed and signed in, run:
\begin{lstlisting}
claude setup-token
\end{lstlisting}
    It opens a browser and then prints a long-lived token. (With an API key, create one in the
    Anthropic console instead.)
  \item In the \ui{Claude} card (figure~\ref{fig:settings-claude}) choose \ui{OAuth token
    (setup-token)} or \ui{Anthropic API key}, paste the value into the field and choose \ui{Save}.
  \item The card's status changes to \ui{OAuth token registered} (or \ui{API key registered}).
    \ui{Replace} and \ui{Remove} change or delete it later.
\end{enumerate}

\screenshot[0.8\linewidth]{25-settings-claude-desktop}{The Claude card: how to get a token, and where to paste it.}{settings-claude}

\section{Your GitHub token}\label{sec:github-token}

The token decides which repositories you see, and your agents push branches, open and merge
pull requests and comment with it. GitHub enforces the access; Agern never widens it.

\begin{enumerate}
  \item On GitHub, under \emph{Settings $\rightarrow$ Developer settings $\rightarrow$ Personal access
    tokens}, create either
    \begin{itemize}
      \item a \textbf{fine-grained token} for the organization with \textbf{Contents},
        \textbf{Issues} and \textbf{Pull requests}: read \& write, plus \textbf{Metadata}: read; or
      \item a \textbf{classic token} with the \texttt{repo} scope.
    \end{itemize}
    These are the permissions the \ui{GitHub} card lists (figure~\ref{fig:settings-github}).
  \item Paste it into \ui{Personal access token} in the \ui{GitHub} card and choose \ui{Save}.
    Agern checks it with GitHub (\ui{Checking with GitHub\ldots}) and refuses a token GitHub rejects.
  \item The card now says \ui{connected as @}\emph{your-login} and shows the kind of token and,
    for a classic token, its scopes. If a classic token lacks \texttt{repo}, a warning asks you to
    create one with the \texttt{repo} and \texttt{read:org} scopes.
\end{enumerate}

\begin{note}[More permissions you may need]
If runs change files under \file{.github/workflows/}, also grant the \texttt{workflow} scope
(classic) or \textbf{Workflows}: read \& write (fine-grained); GitHub refuses such a push
otherwise. An agent can do everything the token allows, so grant it only the repositories Agern
should work on. Scopes, automatic merging and how to restrict it:
\securitydoc{the README's section ``Security: read this first''}.
\end{note}

\screenshot[0.8\linewidth]{26-settings-github-desktop}{The GitHub card with a registered classic token and its scopes.}{settings-github}

\section{Models and cost}\label{sec:models}

Every chat and run you start runs on your Claude account. The \ui{Models \& cost} card
(figure~\ref{fig:settings-models}) lets you choose cheaper models per job and shows your usage.

\begin{enumerate}
  \item Look at the \ui{Usage} block: what you spent through Agern today and this month, and how
    much of Claude's usage limits is used. Claude reports these while one of your agents runs.
  \item To save money, choose the preset \ui{Economy}: Sonnet for chat and spec, Haiku for the four
    reviewers, at most one fix round. \ui{Default} returns to the factory defaults.
  \item Or set a model per role: type an alias (\texttt{opus}, \texttt{sonnet}, \texttt{haiku}) or
    a full model id into the row of \texttt{chat}, \texttt{spec-writer}, \texttt{builder}, a
    reviewer or \texttt{approver}. Empty means the default. \ui{fix rounds} sets how often the
    builder may answer reviewer changes before a human decides.
  \item Choose \ui{Save models}. New chats and runs use the choice; a run keeps the models it
    started with.
\end{enumerate}

\screenshot[0.8\linewidth]{28-settings-models-desktop}{Models \& cost: presets, the usage block and a model per role.}{settings-models}

\section{Optional: image generation}\label{sec:images}

When a request or an issue asks for an icon, an illustration or other graphics, agents can have
them made by an image AI on your account.

\begin{enumerate}
  \item In the \ui{Image generation} card (figure~\ref{fig:settings-images}) choose a provider:
    OpenAI or Google Gemini with your own API key and a model, or OpenAI Codex signed in with your
    ChatGPT account.
  \item For an API-key provider, paste the key, keep or change the model and choose \ui{Save}; the
    provider checks the key first. For Codex, choose \ui{Connect ChatGPT}, open the address shown
    and enter the code.
\end{enumerate}

Without a provider, agents simply do not generate images. Codex is only offered when the server
has the \texttt{codex} tool installed.

\screenshot[0.8\linewidth]{27-settings-images-desktop}{The image generation card with a provider set up.}{settings-images}

\section{Check that it works}\label{sec:check-accounts}

\begin{enumerate}
  \item Choose \ui{Projects} in the top bar. The repositories of the organization that your token
    can read appear as cards (figure~\ref{fig:projects} in the next chapter). If the page says
    \ui{No repositories visible}, your token cannot read any repository of the organization.
  \item Open a project, choose \ui{Start chat} if there is no chat yet, and ask something, for
    example ``What does this repository do?''. The agent answers within seconds, showing the
    tools it uses (figure~\ref{fig:check-chat}).
\end{enumerate}

\screenshot{09-project-streaming-desktop}{A chat answering: the agent reads the code and replies.}{check-chat}

\whatyouhave
Your Claude credential and GitHub token are registered, you have chosen your models, and a chat
answers. Agern is ready for real work. If your projects need a toolchain or a database for their
builds and tests, set up sandboxes in chapter~\ref{ch:sandboxes}; otherwise continue with your
first project in chapter~\ref{ch:project}.
